\documentclass[10pt,xcolor=svgnames]{beamer} %Beamer
\usepackage[justify]{ragged2e}
\usepackage{etoolbox}
\usepackage{lipsum}
\usepackage{natbib}
\usepackage{url}
\usepackage{amsmath}
\usepackage{amsthm}
\usepackage{amsfonts}
\usepackage{mathtools}
\usepackage{amssymb}
\usepackage{gensymb}
\usepackage{amsfonts}
\usepackage{bbm}
\usepackage{caption}
%\usepackage{eufrak}
\usepackage{gensymb}
\usepackage{tikz-cd}
\usepackage{mathdots}
\usepackage{graphicx}
\usepackage{cancel}
\usepackage{multimedia}
\usepackage{hyperref}
\usepackage[T1]{fontenc}
\usepackage[utf8]{inputenc}


\newcommand{\setm}{\smallsetminus}
\newcommand{\N}{\mathbb{N}}
\newcommand{\Z}{\mathbb{Z}}
\newcommand{\Q}{\mathbb{Q}}
\newcommand{\R}{\mathbb{R}}
\newcommand{\C}{\mathbb{C}}
\renewcommand{\H}{\mathbb{H}}
\renewcommand{\P}{\mathbb{P}}
\newcommand{\RP}{\mathbb{RP}}
\newcommand{\CP}{\mathbb{CP}}
\newcommand{\HP}{\mathbb{HP}}
%\newcommand{\G}{\mathbb{G}}
\newcommand{\G}{\mathcal{G}}
\newcommand{\B}{\mathcal{B}}
\newcommand{\DR}{\mathcal{DR}}


\renewcommand{\phi}{\varphi}
\newcommand{\rad}[1]{\sqrt{#1}}
\newcommand{\norm}[1]{\left\lVert#1\right\rVert}
\newcommand{\Qmatrix}[4]{\begin{pmatrix}#1&#2\\#3&#4\end{pmatrix}}
\newcommand{\qmatrix}[4]{\left(\begin{smallmatrix}#1&#2\\#3&#4\end{smallmatrix}\right)}

\newcommand{\g}{\mathfrak{g}}
\newcommand{\h}{\mathfrak{h}}
\DeclareMathAlphabet{\mathpzc}{OT1}{pzc}{m}{it} % k
\renewcommand{\k}{\mathpzc{k}}
\newcommand{\q}{\mathfrak{q}}
%\renewcommand{\c}{\mathfrak{c}}
\newcommand{\gl}{\mathfrak{gl}}
\renewcommand{\sl}{\mathfrak{sl}}
\newcommand{\so}{\mathfrak{so}}
\renewcommand{\u}{\mathfrak{u}}
\newcommand{\su}{\mathfrak{su}}
\renewcommand{\sp}{\mathfrak{sp}}
\newcommand{\m}{\mathfrak{m}}
\newcommand{\aff}{\mathfrak{aff}}
\usepackage{mathrsfs}
\newcommand{\F}{\mathscr{F}}
\newcommand{\X}{\mathfrak{X}}
\renewcommand{\L}{\mathscr{L}}
\renewcommand{\O}{\mathscr{O}}
\newcommand{\M}{\mathcal{M}}

\newcommand{\1}{\mathbbm{1}}
\newcommand{\J}{\mathbb{J}}
\newcommand{\A}{\mathbb{A}}
\newcommand{\z}{\mathfrak{z}}
\DeclareMathOperator{\dd}{d}
\DeclareMathOperator{\dx}{dx}
\DeclareMathOperator{\Diff}{Diff}
\DeclareMathOperator{\Isom}{Isom}
\DeclareMathOperator{\Ad}{Ad}
\DeclareMathOperator{\ad}{ad}
\DeclareMathOperator{\Aut}{Aut}
\DeclareMathOperator{\aut}{aut}
\DeclareMathOperator{\End}{End}
\DeclareMathOperator{\sym}{Sym}
\DeclareMathOperator{\tr}{tr}
\DeclareMathOperator{\der}{der}
\DeclareMathOperator{\Mor}{Mor}
\DeclareMathOperator{\Map}{Map}
\DeclareMathOperator{\Ob}{Ob}
\DeclareMathOperator{\Hom}{Hom}
\apptocmd{\frame}{}{\justifying}{} % Allow optional arguments after frame.
%\usepackage{quiver}
\usepackage{palatino, dsfont} %font type
\usefonttheme{metropolis} %Type of slides
\usefonttheme[onlymath]{serif} %font type Mathematical expressions
\usetheme[progressbar=frametitle,titleformat frame=smallcaps,numbering=counter]{metropolis} %This adds a bar at the beginning of each section.
\useoutertheme[subsection=false]{miniframes} %Circles in the top of each frame, showing the slide of each section you are at

\usepackage{appendixnumberbeamer} %enumerate each slide without counting the appendix
\setbeamercolor{progress bar}{fg=Maroon!70!Coral} %These are the colours of the progress bar. Notice that the names used are the svgnames
\setbeamercolor{title separator}{fg=DarkSalmon} %This is the line colour in the title slide
\setbeamercolor{structure}{fg=black} %Colour of the text of structure, numbers, items, blah. Not the big text.
\setbeamercolor{normal text}{fg=black!87} %Colour of normal text
\setbeamercolor{alerted text}{fg=DarkRed!60!Gainsboro} %Color of the alert box
\setbeamercolor{example text}{fg=Maroon!70!Coral} %Colour of the Example block text


\usepackage{media9}%
\newcommand{\includemovie}[3]{%
	\includemedia[%
	width=#1,height=#2,%
	activate=pagevisible,%
	deactivate=pageclose,%
	addresource=#3,%
	flashvars={%
		src=#3 % same path as in addresource!
		&autoPlay=true % default: false; if =true, automatically starts playback after activation (see option ‘activation)’
		&loop=true % if loop=true, media is played in a loop
		&controlBarAutoHideTimeout=0 %  time span before auto-hide
	}%
	]{}{StrobeMediaPlayback.swf}%
}% end of the new command



\setbeamercolor{palette primary}{bg=NavyBlue!50!DarkOliveGreen, fg=white} %These are the colours of the background. Being this the main combination and so one. 
\setbeamercolor{palette secondary}{bg=NavyBlue!50!DarkOliveGreen, fg=white}
\setbeamercolor{palette tertiary}{bg=NavyBlue!40!Black, fg= white}
\setbeamercolor{section in toc}{fg=NavyBlue!40!Black} %Color of the text in the table of contents (toc)

%These next packages are the useful for Physics in general, you can add the extras here. 
\usepackage{amsmath,amssymb}
\usepackage{slashed}
\usepackage{cite}
\usepackage{relsize}
\usepackage{caption}
\usepackage{subcaption}
\usepackage{multicol}
\usepackage{booktabs}
\usepackage[scale=2]{ccicons}
\usepackage{pgfplots}
\usepgfplotslibrary{dateplot}
\usepackage{geometry}
\usepackage{xspace}
\newcommand{\themename}{\textbf{\textsc{bluetemp}\xspace}}%metropolis}}\xspace}

\title{Compactification of the fifth Painlevé Foliation}
\author[Name]{Mattia Morbello} %With inst, you can change the institution they belong
\subtitle{Web-seminar on Painlevé Equations and related topics}
\institute[uni]{Université de Rennes}
\date{8 April 2026} %Here you can change the date %You can modify the location of the logo by changing the command \vspace{}. 




\begin{document}
{
\setbeamercolor{background canvas}{bg=NavyBlue!50!DarkOliveGreen, fg=white}
\setbeamercolor{normal text}{fg=white}
\maketitle
%\includegraphics[width=3cm]{LOGO.jpg}
}%This is the colour of the first slide. bg= background and fg=foreground







\section{Meromorphic Connections}


\begin{frame}{Meromorphic Connections on $\P^1$}
\metroset{block=fill}
\begin{exampleblock}{Definition}
A rank 2 meromorphic connection $(E, D, \nabla)$ on $\P^1$ is the data of:
\begin{itemize}
	\item a rank 2 holomorphic vector bundle $E\to \P^1$,
	\item a effective divisor $D$ of $\P^1$ called the polar divisor,
	\item a morphism $\nabla\colon\mathcal E\to \mathcal E\otimes \Omega^1(D)$.
\end{itemize}
\end{exampleblock}
\begin{exampleblock}{Remark}
	The sheaf $\Omega^1(D):=\Omega^1\otimes \O(D)$ is the sheaf of meromorphic 1-forms with poles as prescribed by $D$.
\end{exampleblock}
\begin{exampleblock}{Leibniz Rule}
The operator $\nabla$ is required to satisfy
\[\nabla(f\sigma)=df\cdot\sigma+f\nabla\sigma\]
\end{exampleblock}
\end{frame}



%\begin{frame}{An Example}
%\metroset{block=fill}
%\begin{exampleblock}{Euler System: $X=\P^1$ and $\deg D=2$}
%$(\nabla, E, D)=(d+\Omega,\; \O\oplus\O,\; [0]+[\infty])$.
%\[d\begin{pmatrix}y_1\\y_2\end{pmatrix}+\Omega\begin{pmatrix}y_1\\y_2\end{pmatrix}\;\;\;=\;\;\;\begin{pmatrix}dy_1\\dy_2\end{pmatrix}+\Bigg[\underbrace{\begin{pmatrix}a&b\\c&d\end{pmatrix}\frac{dx}{x}}_{\Omega}\Bigg]\begin{pmatrix}y_1\\y_2\end{pmatrix}\]
%\begin{itemize}
%	\item Has simple poles in $x=0$ and $x=\infty$.
%	\item Has a solution matrix of the form \[Y(x)=x^{A}:=\exp\Big(A\log x\Big)\]
%	\item Has non-trivial \textit{monodromy} due to the complex logarithm.
	%\item It is a \textit{normal form}.
%	\item Any deformation of parameters modify the monodromy (rigidity).
%\end{itemize}
%\end{exampleblock}
%\end{frame}


%\begin{frame}{An example}
%	\metroset{block=fill}
%	\begin{exampleblock}{Meromorphic connections on the trivial bundle}
%		Let $E\cong\O\oplus\O$ be the trivial bundle. Any $(E, D, \nabla)$ is of the form 
%		\[\nabla=d+\Omega, \;\;\;\text{ for }\;\;\;\Omega\in \mathfrak{gl}_2\left(\Omega^1(D)\right).\]
%		The action on a global section is then the following:
%		\[\nabla Y=\nabla\begin{pmatrix}y_1\\y_2\end{pmatrix}=\begin{pmatrix}dy_1\\dy_2\end{pmatrix}+\begin{pmatrix}\omega_{1,1}&\omega_{1,2}\\\omega_{2,1}&\omega_{2,2}\end{pmatrix}\begin{pmatrix}y_1\\y_2\end{pmatrix}\]
%		for some $\omega_{i,j}\in \Omega^1(D)$
%	\end{exampleblock}
%\end{frame}


%\begin{frame}{What if $E$ is not trivial ?}
%\metroset{block=fill}
%\begin{exampleblock}{Remark (Birkhoff-Grothendieck Theorem)}
%	Any rank 2 vector bundle $E$ on $\P^1$ is of the form 
%	\[E\cong\O(k_1)\oplus\O(k_2), \;\;\;\text{ for }\;\;\;k_1,k_2\in \Z.\]
%\end{exampleblock}
%\vskip 15pt
%\begin{exampleblock}{Remark}
%	The cover $\{U_0,U_\infty\}$ trivialize any vector bundle $E$ on $\P^1$. 
%\end{exampleblock}
%\end{frame}




\begin{frame}{Local Expression}
	\metroset{block=fill}
	\begin{exampleblock}{Trivialisation}
		On a trivializing open set $U$, a connection $(E,D,\nabla)$ has the form
		\[\nabla_{|U}=d+\Omega_U, \;\;\text{ for }\;\;\Omega_U\in \mathfrak{gl}_2(\Omega^1(D\cap U)).\]
	\end{exampleblock}
	\begin{exampleblock}{Gluing conditions}
		Any connection $(E, D, \nabla)$ on $\P^1$ is then identified with the data
		\[\begin{cases*}d+\Omega_0 \;\;\text{ on }U_0:=\P^1\setminus\{\infty\}\\d+\Omega_\infty \;\text{ on }U_\infty:=\P^1\setminus\{0\}\\g_{0,\infty}\;\;\;\;\;\;\text{ cocycle of } E\end{cases*}\]
		with gluing condition on the overlap
		\[\Omega_0=g_{0,\infty}^{-1}\cdot \Omega_\infty\cdot g_{0,\infty}+ g_{0,\infty}^{-1}\cdot dg_{0,\infty}.\]
	\end{exampleblock}
\end{frame}





%\begin{frame}{Local Monodromy}
%	\begin{exampleblock}{Horizontal sections}
%		A local section $Y$ is said horizontal if $\nabla_{|U} Y=dY+\Omega_UY=0$
%	\end{exampleblock}
%	\begin{exampleblock}{Monodromy}
%	\begin{center}
%	\end{center}
%	\[Y^\gamma=M_\gamma\cdot Y\;\;\text{ for }\;\;M_\gamma\in \mathrm{GL}_2(\C).\]
%	\end{exampleblock}
%\end{frame}










\begin{frame}{Equivalence of Connections}
\metroset{block=fill}
\begin{exampleblock}{Gauge equivalence of connections}
We say that $(E, D, \nabla)\sim(E',D',\nabla')$ if there exists $\Phi\in\mathrm{Mor}(E, E')$ such that 
\[\Phi^*\nabla'=\nabla.\]
\end{exampleblock}
\begin{exampleblock}{Fact}
	\[(E, D, \nabla)\sim(E',D',\nabla')\;\;\;\implies\;\;\;M_\nabla\sim M_{\nabla'}\]
\end{exampleblock}
\begin{exampleblock}{Local equivalence}
The connections matrices on a trivialising open set $U$ are locally related by: 
\[\Omega_U=\Phi^{-1}\Omega_U'\Phi+\Phi^{-1}d\Phi\]
\end{exampleblock}
\end{frame}



\begin{frame}{Equivalence of Connections}
	\metroset{block=fill}
	\begin{exampleblock}{Example}
	The action of the meromorphic gauge transformation
	\[\Phi(x)=\begin{pmatrix}1&0\\0&\frac{1}{x-q}\end{pmatrix}\]
	on the connection $(\O\oplus\O, D, d+\Omega)$ is
	\[(\O\oplus\O(1), D+[q], d+\Omega'),\]
	where
	\[\Omega':=\begin{pmatrix}\omega_{1,1}&(x-q)\omega_{1,2}\\\frac{1}{x-q}\omega_{2,1}&\omega_{2,2}-\frac{dx}{x-q} \end{pmatrix}\]
	\end{exampleblock}
\end{frame}



%\begin{frame}{Meromorphic Gauge Transformations}
%\metroset{block=fill}
%\begin{exampleblock}{Remark}
%Meromorphic gauge transformations do not change the monodromy and changes the degree of the bundle.
%\end{exampleblock}
%\vfill
%\begin{exampleblock}{Example}
%\[\Psi:=\begin{pmatrix}1&0\\0&\frac{1}{x-q}\end{pmatrix}\]
%is a meromorphic gauge transformation, then:
%\[\Psi\;\colon\; (\nabla, \O\oplus\O, D) \mapsto (\tilde\nabla, \O\oplus\O(1), D+[q]),\]
%and the pole $q$ has trivial monodromy.
%\end{exampleblock}
%\begin{exampleblock}{Conclusion}
%Since we are interested in isomonodromic deformations, we have to study connections up to meromorphic gauge transformations.
%\end{exampleblock}
%\end{frame}




\section{Painlevé V Connections}





%\begin{frame}{Painlevé Equations}
%\metroset{block=fill}
%\begin{exampleblock}{Isomonodromic Deformations}
%If $\deg D \geq 4$, we can deform $[\nabla]$ keeping the (conjugacy class of the) monodromy constant.
%\end{exampleblock}
%\vfill
%set image of coalescent diagram ?
%\begin{exampleblock}{Painlevé Equations (Jimbo, Miwa, Ueno 1981)}
%Are second order differential equations whose solutions parametrise isomonodromic deformations of rank 2 meromorphic connections with $\deg D = 4$.
%\end{exampleblock}
%\end{frame}










\begin{frame}{Connexions du type PV}
	\metroset{block=fill}
	\begin{exampleblock}{Definition}
		Meromorphic connections $(E, D, \nabla)$ such that
		\[D^{\text{min}}=[0]+2[1]+[\infty],\]
		where $D^{\text{min}}$ is the \textit{minimal} polar divisor w.r.t. the gauge equivalence class of  $(E, D, \nabla)$.
	\end{exampleblock}
	\vfill
	\begin{exampleblock}{Consequence}
		In $D^{\text{min}}$ only appear poles
		\begin{itemize}
			\item that cannot be eliminated nor reduced via a gauge transformation,
			\item with minimal Poincaré rank,
			\item with non-trivial monodromy.
		\end{itemize}
	\end{exampleblock}
\end{frame}












\begin{frame}{Normal Form}
\metroset{block=fill}
\begin{exampleblock}{Normal Form on $\O\oplus\O(2)$ (Diarra, Loray 2019)}
\begin{align*}
	\nabla_{|0}= & \,d+\Omega_0=\\\\&d+\begin{pmatrix}0&1\\0&t\end{pmatrix}\frac{dx}{(x-1)^2}+\begin{pmatrix}0&-1\\0&\kappa_1\end{pmatrix}\frac{dx}{x-1}+\begin{pmatrix}0&1\\0&-\kappa_0\end{pmatrix}\frac{dx}{x}\\\\&+\begin{pmatrix}0&0\\\kappa_\infty&0\end{pmatrix}xdx+\begin{pmatrix}0&0\\p&-1\end{pmatrix}\frac{dx}{x-q}+\begin{pmatrix}0&0\\\hat K&0\end{pmatrix}dx,
\end{align*}
\end{exampleblock}
\begin{exampleblock}{Fixed Parameters}
\[\Theta:=\{\kappa_0, \kappa_1, \kappa_\infty\}\subseteq \C\]
\end{exampleblock}
\end{frame}





\begin{frame}{Geometrical Interpretation of $t$ and $q$}
	\metroset{block=fill}
	\begin{exampleblock}{Action of $f\in\mathrm{Aut}(\P^1)$}
		Let $f\in\mathrm{Aut}(\P^1)$, then:
		\begin{itemize}
			\item $q\mapsto f(q)$,
			\item $t\mapsto Df(1)\cdot t$.
		\end{itemize}
	\end{exampleblock}
	\begin{center}
		\includegraphics[width=5cm]{ratirr.jpeg}
	\end{center}
\end{frame}


\begin{frame}{Rational Irregular Curves}
	\metroset{block=fill}
	\begin{exampleblock}{Definition}
		A rational irregular curve is $(\P^1, D, J)$, where 
		\begin{itemize}
			\item $D=\sum_{i\in I} n_i[p_i]$ is an effective divisor,
			\item $J=\{j^{n_i-1}(p_i)\}_{i\in I}$ is a collection of jets.
		\end{itemize}
	\end{exampleblock}
	\begin{exampleblock}{Example}
		$(\P^1, D, J)$ with 
		\begin{itemize}
			\item $D=[0]+2[1]+[\infty]+[q]$
			\item $J=\{\;j^0(0),\; j^1(1)=(1,t),\; j^0(\infty),\; j^0(q)\;\}$.
		\end{itemize}
	\end{exampleblock}
	\begin{exampleblock}{Proposition}
		$[(\nabla, E, D)]$ "generic" of PV type $\iff$ a rational irregular curve and $p\in\C$.
	\end{exampleblock}
\end{frame}




%\begin{frame}{Geometrical Interpretation of $p$}
%\metroset{block=fill}
%\begin{exampleblock}{Projectivisation of the Connection}
%Given a Painlevé V connection $(\nabla, \O\oplus\O(2))$, we can consider
%\begin{itemize}
%	\item $\P(E)$, that is a $\P^1$-bundle with fiber coordinate $y:=y_2/y_1$
%	\item The Riccati differential equation \[dy=-\omega_{1,2}y^2+\omega_{2,2}y+\omega_{2,1}\] induced by $\nabla$.
%\end{itemize}
%\end{exampleblock}
%\begin{exampleblock}{Riccati Foliation}
%\begin{itemize}
%	\item $\mathbb{F}_2:=\P(\O\oplus\O(2))$ the second Hirzebruch surface
%	\item $\mathcal F$ the foliation on $\mathbb{F}_2$ induced by the differential equation below.
%\end{itemize}
%\end{exampleblock}
%\end{frame}







%\begin{frame}{Geometrical Interpretation of $p$}
%\metroset{block=fill}
%\begin{center}
%    \includegraphics[width=11cm]{DISEGNI(16).jpeg}
%\end{center}
%\end{frame}



%\begin{frame}{Monodromy}
%\metroset{block=fill}
%\begin{exampleblock}{Fact}
%\[\mathrm{Mon}(\nabla, \O\oplus\O(2))=\mathrm{Mon}(\mathcal{F}, \mathbb{F}_2)\]
%\end{exampleblock}
%\vfill
%\begin{exampleblock}{Stokes Phenomena}
%The local monodromy around $x=1$ is more complicated: solutions are defined only on sectors around the singularity and the change of sector give rise to a monodromy.
%\end{exampleblock}
%\end{frame}



%\begin{frame}{Monodromy}
%\metroset{block=fill}
%\begin{exampleblock}{Problem}
%Isomonodromic deformations
%\end{exampleblock}
%\begin{exampleblock}{Tools}
%begin{itemize}%
%	\item Moduli space of connections
%	\item Deligne-Mumford compactification via stable curves
%	\item Painlevé V differential equation
%\end{itemize}
%end{exampleblock}
%vfill
%\begin{exampleblock}{Painlevé Equations (Jimbo, Miwa, Ueno 1981)}
%Are second order differential equations whose solutions parametrise isomonodromic deformations of rank 2 meromorphic connections with $\deg D = 4$.
%\end{exampleblock}
%\end{frame}



%\begin{frame}{Questions}
%	\metroset{block=fill}
%	\begin{exampleblock}{Isomonodromic deformations}
%		How can we deform the parameters $(t,q,p)$ without changing the monodromy of the connection?\\
%		What happens when $t=0$ or $t=\infty$?
%	\end{exampleblock}
%	\vfill
%	\begin{exampleblock}{Tools}
%		\begin{itemize}
%			\item Moduli space of PV connections,
%			\item Painlevé V equation,
%			\item Isomonodromic foliation.
%		\end{itemize}
%	\end{exampleblock}
%\end{frame}



\section{Moduli Space and Compactification}



\begin{frame}{Moduli Space}
	\metroset{block=fill}
	\begin{exampleblock}{Proposition}
		$[(\nabla, E, D)]$ "generic" of PV type $\iff \begin{cases} (\P^1, D, J)\;\;,\text{ and}\\p\in \C\end{cases}$
	\end{exampleblock}
	\begin{exampleblock}{Moduli space of rational irregular curves $\mathcal M$}
		\begin{align*}
			\mathcal M:=\Big\{q\in\P^1\setminus\{0,1,\infty\};\;\;t\in T_{1}\P^1\setminus \{0\}\Big\}.
		\end{align*}
	\end{exampleblock}
	\begin{exampleblock}{Moduli space of PV connections}
		\[\mathfrak{Con}^V_\Theta\supseteq\mathcal M\times \C\]
	\end{exampleblock}
\end{frame}



\begin{frame}{Moduli Space $\M\subseteq \P^1\times\P^1$}
	\metroset{block=fill}
	\begin{center}
		\includegraphics[width=9cm]{Espmod.jpeg}
	\end{center}
	\begin{exampleblock}{Remark}
		The moduli space $\M$ is a generalization of $\M_{0,5}$ where the fifth marked point is replaced by a tangent vector.
	\end{exampleblock}
\end{frame}





\begin{frame}{Moduli Space of Conics}
	\metroset{block=fill}
	\begin{center}
		\includegraphics[width=10cm]{image9.jpg}
	\end{center}
	\begin{exampleblock}{Moduli Space of Smooth Conics}
		\[\P^2\setminus\Big(\overline{AB}, \overline{AC}, \overline{BC}, l \Big)\cong \M\]
	\end{exampleblock}
\end{frame}



\begin{frame}{Blow-up of $A$}
	\metroset{block=fill}
	\begin{center}
		\includegraphics[width=7cm]{image10.png}
	\end{center}
\end{frame}



\begin{frame}{Blow-up of $B$, $C$ and $L$}
	\metroset{block=fill}
	\begin{center}
		\includegraphics[width=8cm]{image11.png}
	\end{center}
\end{frame}




\begin{frame}{Compactification $\overline{\mathcal M}$}%
	\metroset{block=fill}
	\begin{center}
		\includegraphics[width=9cm]{image12.png}
	\end{center}
\end{frame}





\begin{frame}{Compactification $\overline{\mathcal M}$}
\metroset{block=fill}
\begin{exampleblock}{Compactification (M., 2025)}
Following, \textit{mutatis mutandis}, a similar idea that Kapranov used to compactify the spaces $\M_{0,n}$, we get:
\begin{center}
	\includegraphics[width=11cm]{image0.jpeg}
\end{center}
where $\widetilde\M$ is the weak del Pezzo surface of degree five.
\end{exampleblock}
\end{frame}



\begin{frame}{Universal Curve}
	\metroset{block=fill}
	\begin{center}
		\includegraphics[width=8cm]{image1-5.jpeg}
	\end{center}
\end{frame}



\begin{frame}{Compactification $\overline{\mathfrak{Con}_V^\Theta}$}
	\metroset{block=fill}
	\begin{exampleblock}{Line Bundle Extension (M., 2025)}
		The trivial bundle $\mathcal M\times \C$ extends to
		\[\O_{\widetilde{\mathcal M}}\Big(A_0+2Q^1+2B_0^1-B_\infty^0-B_\infty^\infty\Big).\]
		Moreover, it is trivial when restricted to $Q^0$, $Q^\infty$ and $Q^1\cup A_\infty$.
	\end{exampleblock}
	\begin{center}
		\includegraphics[width=5cm]{DISEGNI (16).jpeg}
	\end{center}
\end{frame}



\begin{frame}{Compactification $\overline{\mathfrak{Con}_V^\Theta}$}
	\begin{center}
		\includegraphics[width=7cm]{CPXGEOM2.png}
	\end{center}
\end{frame}



\begin{frame}{Compactification $\overline{\mathfrak{Con}_V^\Theta}$}
	\metroset{block=fill}
	\begin{exampleblock}{Theorem}
		The compactified moduli space $\overline{\mathfrak{Con}_V^\Theta}$ comes with a fibration $\pi\colon\overline{\mathfrak{Con}_V^\Theta}\xrightarrow{t}\P^1$ such that:
		\begin{itemize}
			\item For any $t\in\C^*$, the fiber $\pi^{-1}(t)$ is an Okamoto space, that is a 8-blow-up of the second Hirzebruch surface $\mathbb F_2$.
		\end{itemize}
	\end{exampleblock}
	\begin{center}
	\includegraphics[width=5cm]{okam.png}
	\end{center}
\end{frame}




\begin{frame}{Compactification $\overline{\mathfrak{Con}_V^\Theta}$}
	\metroset{block=fill}
	\begin{exampleblock}{Theorem (M, 2025)}
		The compactified moduli space $\overline{\mathfrak{Con}_V^\Theta}$ comes with a fibration $\pi\colon\overline{\mathfrak{Con}_V^\Theta}\xrightarrow{t}\P^1$ such that:
		\begin{itemize}
			\item The surface $\pi^{-1}(0)$ is given by $\mathcal A_0\cup\mathcal B_0^1$. Both these components are isomorphic to $\mathbb F_1$.
			\item Let us denote by $\mathcal F^+_\infty$ and $\mathcal F^-_\infty$ the blow up of the two special sections in $\mathcal A_\infty$. The surface $\pi^{-1}(\infty)$ is given by $\mathcal B_0^\infty\cup\mathcal B_\infty^\infty\cup\mathcal F^+_\infty\cup\mathcal F^-_\infty$. The surfaces $\mathcal B_0^\infty$ and $\mathcal B_\infty^\infty$ are isomorphic to $\mathbb F_1$.
		\end{itemize}
	\end{exampleblock}
\end{frame}





\section{Painlevé V Foliation}


%\begin{frame}{Painlevé V Equation}
%\metroset{block=fill}
%\begin{exampleblock}{The Equation (PV)}
%\begin{align*}
%    q''(t)=&\Bigg(\frac{1}{2q(t)}+\frac{1}{q(t)-1}\Bigg)q'(t)^2-\frac1tq(t)'\\\\&+\frac{(q(t)-1)^2}{t^2}\Bigg(\alpha q(t)+\frac{\beta}{q(t)}\Bigg)+\gamma \frac {q(t)}{t}+\delta\frac{q(t)(q(t)+1)}{q(t)-1}
%\end{align*}
%Where $\alpha, \beta, \gamma, \delta$ are parameters depending on $\Theta$.
%\end{exampleblock}
%\end{frame}


\begin{frame}{Monodromy of a PV Connection}
	\metroset{block=fill}
	\begin{exampleblock}{Wild fundamental groupoid (E. Paul, J-P. Ramis, 2023)}
		\begin{center}
			\includegraphics[width = 11cm]{wildfund.png}
		\end{center}
	\end{exampleblock}
\end{frame}

\begin{frame}{Monodromy of a PV Connection}
	\metroset{block=fill}
	\begin{center}
		\includegraphics[width = 5cm]{wildfund2.png}
	\end{center}
	\begin{exampleblock}{Fact}
		The monodromy along $\gamma_{1,2}^+\star\gamma_{1,2}^-$ is given by
		\[M^\gamma=S_0S_1D\]
		with $S_0$ upper triangular, $S_1$ lower triangular (Stokes matrices) and $D$ diagonal.
	\end{exampleblock}
\end{frame}




\begin{frame}{Painlevé V Equation}
\metroset{block=fill}
\begin{exampleblock}{The Hamiltonian Function}
	\[H^V:=\frac{q \left(q -1\right)^{2} p^{2}-\left(\kappa_0 \left(q -1\right)^{2}+\kappa_1  q \left(q -1\right)-t q \right) p +\kappa_\infty \left(q -1\right)}{t}\]
\end{exampleblock}
\begin{exampleblock}{The Hamiltonian System}
\[\begin{cases}
    \displaystyle\frac{\partial H^V}{\partial p}=\displaystyle\frac{dq}{dt}\\ \\\displaystyle\frac{\partial H^V}{\partial q}=-\displaystyle\frac{dp}{dt}
\end{cases}\]
\end{exampleblock}
\end{frame}



\begin{frame}{How to Find First Integrals}
	\metroset{block=fill}
	\begin{exampleblock}{Deformation of a PV connection}
		A PV connection splits into a connection over any smooth component. Both connections have a new pole at the node.
	\end{exampleblock}
	\begin{center}
		\includegraphics[width=10cm]{stableconn.png}
	\end{center}
\end{frame}




\begin{frame}{How to Find First Integrals}
	\metroset{block=fill}
	\begin{exampleblock}{Proposition}
		The residual spectral datum in the node is a first integral for the hamiltonian vector field in restriction to the respective boundary component.
	\end{exampleblock}
	\begin{exampleblock}{Idea}
	Since we are following an isomonodromic leaf, the monodromy along the loop $\gamma$ must be the same as the one of $\gamma_1$ and $\gamma_2$, corresponding to the monodromy of the new poles at the node. In particular the residual spectral datum in the node must be preserved during the deformation.
	\end{exampleblock}
\end{frame}




\begin{frame}{Isomonodromic Foliation On the Okamoto Divisor}
	\metroset{block=fill}
	\begin{exampleblock}{Proposition}
		The isomonodromic foliation is stationary along the Okamoto divisor.
	\end{exampleblock}
	\begin{center}
		\includegraphics[width=6cm]{folQ1.png}
	\end{center}
\end{frame}




\begin{frame}{First Integrals around $t=0$}
\metroset{block=fill}
\begin{center}
    \includegraphics[width=8cm]{FI0.png}
\end{center}
\end{frame}








\begin{frame}{First Integrals around $t=\infty$}
\metroset{block=fill}
\begin{center}
    \includegraphics[width=10cm]{FIinf.png}
\end{center}
\end{frame}






\begin{frame}{Thanks :)}
\begin{center}
	\Large Thanks for your attention !!
\end{center}
\end{frame}





\end{document}